% ================================================================
%  18.7 Les fruits du soir (原书 p.908–916)
%  源文: _work/chunks/chunk14.md
% ================================================================
\section{(7) 暮熟的果实}

\subsection{(1) 尊重}   % 18.7.1

\notehead{179}{}

(4月2日)已经五个星期了(自2月26日那道“沉默”之注起——它揭开了集结在“四种运算”名目
之下的那组笔记),我一直在重新巡视有关\emph{葬礼}的种种“物质性”的、或者(多少)“技术性”
的事实。在“四种运算”里,我把自己限定在严格意义下的“诈骗”这一面——也就是越过了那道
“门槛”的一面;那道门槛,同名的注(n\textsuperscript{o}~172)里谈过,它分开\textbf{恶劣的心性}(表现为
“自动排斥”的反射,往往不顾一个数学家最起码的直觉)与公然的\textbf{蓄意刁难}、与铁板钉钉的剽
窃。在刚写完的“荒废的工地”里,我面对的主要是\emph{葬礼}的“第一层次”,门槛以下的那一面:
一场广阔视野\dmark{}\origpage{1219}与一些强有力执念观念的埋葬——当然,没有谁非得接着做它们不可,人人都有权
无视或遗忘它们——只是这样一来,也就“把自己埋葬了”,让自己的工作(至少让那部分直接蒙受
被拒视野的工作)或多或少归于不育。

在那里,我总算有了走完一遍的印象!就“工地巡视”(撂荒的那些)而言,它让我对我著作的\emph{葬
礼}有了一份原原本本的把握,同时又让我同十五年来失了音讯的一些主题,多少重新接上了头。这
尤其让我看清了:在接下来几卷《数学反思》里我打算白纸黑字写下的东西,孰急孰缓。我的宗旨,当
然不再是给一门临盆中的科学打细致的地基——这种事我做得够多了;倘若再没有旁人肯像我当年
那样,把自己交给这样一项任务,那就各人自便,大家都自便!我的宗旨毋宁是:提炼某些执念观
念——它们服务于一个1955到1970年间诞生的总体视野;而今天我重新找到它们(多亏我某些昔日学生的努力,又蒙全体默认)时,它们或被遗忘,或沦为笑柄,或被肆无忌惮地据为己有、
残害得失去大半元气。今天重拾它们,我终于松开了缰绳,放开心中那股求知的冲动——六十年代
里,为了没完没了的“服务”任务,我常常只给它一份可怜的口粮。那个时代一去不返了——然而我
知道,在我数学激情的这个新阶段里,服务的冲动,丝毫不减当年。我依旧要为那个由渴望知道的
头脑结成的理想“共同体”服务\fnote{1014}{(*)}{关于“数学共同体”,我第一次表明看法,是在《收获与播种》第一部分“‘数学共同体’:虚构与现实”一节(n\textsuperscript{o}~10)里。在这里我援引一个“由渴望知道的头脑结成的理想共同体”,看上去似乎我又退回到某种东西上去了——它的虚构性质,在所引那一节里曾显示得清清楚楚。但在《愚妄与更新》第八部分里,我有生以来第一次(迟做总比不做好……)确认了我自己的“求知历险”在数学层面上的一个集体维度。(关于这一点,见“孤独的历险”与“一段过去的重量”两节,n\textsuperscript{o}s~47、50,尤其打字稿第134、135页。)同样清楚的是:经历这场集体历险的那个“共同体”(或“集体”),与任何社会学意义上的实体——它体现于特定“时代”的特定“圈子”,带着特定的“心态”,或者(今天)带着特定的权力结构与阶级利益——完全是两回事。我所援引的这个“理想共同体”,“在空间与时间上都无边界”,对我而言丝毫不比社会学实体更不“真实”。它甚至更其根本,因为正是它(如我在同一句话的后文里所写)“继续给我的数学投入以一种比个人消遣、比自我扩张手段更深的意义”。它并不比我自己更“虚构”——我自觉是它的一分子,而且比当年更清醒。“虚构”不在于觉察到这样一个“共同体”的存在,而在于把它同一个我曾自认同属的圈子混为一谈。},它继
续给我的数学投入,以一种比个人消遣、比自我扩张的手段更深的意义。

\pmark\origpage{1220}在这些投入里,当然,“老板”也并没有比当年更缺席。面对那些在我眼里曾是数学世界里
“我的亲近者”的人的恶意与嘲弄,面对那些我爱过、也曾无条件信任的人对我最起码体面感的屡
屡伤害,我心中涌起一股按捺不住的冲动:在那些丢失了敬意的人面前,\textbf{以对这些活生生的、
茁壮而美丽的东西的敬意,来见证我对自己的敬意}——这些东西,是我亲手带到光天化日之下
的。也许,见证这份敬意的最好方式,就是在这段还归于我的宝贵岁月里,做这些东西的仆
人。因此,未来几年里,我打算在《数学反思》的后续各卷中发展的数学思考,将既是重新拾起一个\textbf{孩
童的游戏},又是一件\textbf{服务的礼物},又是一个\textbf{敬意的举动}。

在给\emph{葬礼}画上句号之前,我还想在“物质事实”之外,就这场思考教给我的东西,作一小结。
我先看它教给我的关于他人的事,最后再谈它教给我的关于我自己的事。

在整个思考过程中大白于天下的种种事实里,至今最触目的,仍然是七十年代、八十年代数学世界
里\textbf{风气与人心的败坏}。这种败坏,借成百上千桩“小节”表现出来,如同过去八九年间,一阵阵
回到我耳边的那些“小节”——它们却足以让人错愕,足以催生《收获与播种》第一部分的思考及
其根本追问:事情是怎么(又是何时)走到这一步的?而在这场我今天目睹的阴险而无情的败坏之
中,我的角色是什么,我的位置又在哪里?

\pmark\origpage{1221}这场败坏,在“SGA~4½--SGA~5”那样的行动里,或者在(更难以置信的)反常层会议那场
行动里,达到了顶点——其犬儒与轻蔑,远远超出了我所能想象的一切;哪怕就在我身不由己地发现
它们的前一天,也想象不到。

在这里,不适合回头去谈那些“小节”(其中不少,已在思考途中顺带指出),也不适合谈那些大行
动(它们自有小伎俩伺候)。显露在两者——小节与大骗局——之中的,是同一种精神。在可容忍与
下作之间有时不妨画一道“门槛”,这道门槛本身就相当脆弱、相当人为,是一根护栏;反正(看样
子)已经没有人在乎它了。经由这场把我本人牵连至深的\emph{葬礼},我得以也许比以往任何时候都
更近地看清这种精神——它当然既不是这场葬礼(为我这卑不足道的人发起)的专利,也不是数学
家世界的专利。对此我能说的只是:我不知道,在我们这个时代之外的任何时代,这种精神曾统治
过那个世界,或其他任何一门科学。这大概是众多征兆之一,表明一种文明——以及那无论如何还
在给它赋予意义的东西——分解到了终局。

\pmark 近来,我的思绪不止一次地流连于一个奇怪的巧合:十五年前我告别数学舞台,是受了科学
界某种腐败的冲击——那种腐败,我曾长期选择闭目不见(自以为离它远着呢)。而我是在那所我
本打算终老于此的机构里,猝然同它打了个照面\fnote{1015}{(*)}{关于这一点,见注“有益的挣脱”(n\textsuperscript{o}~42),以及“兄弟与夫妻——或双重署名”及其子注(n\textsuperscript{o}s~134、134$_{1}$)。}。在那里,事关科学家们与军事机器之间那几近普遍的、有利可图的
串通。军方对整个科学世界的这种阴险把持,同样是个新现象,只是(至少就我们今日所知的规模而
言)自第二次世界大战以来才出现。当然,那场“冲击”之所以搅乱了我预定的轨迹(我自己、人人
都预定的轨迹),以至于触发我无可挽回地离开那个我一直自认同属的世界(除了一项默而不宣的
保留……),是因为我心中有一种紧迫而急切的更新之需——这一点,我是事后回望才意识到的
\origpage{1222}。后来,我曾倾向于把触发那次不寻常离开的特殊诱因说得轻些。然而我也知道,那些惯性之力
有多么巨大(同时又是多么无形):它们要把我们永远拴在同一条“轨迹”上,敌视内在的更新——
这也让我掂量出,把我从一条划得那样牢固的轨迹上拔出来,需要何等猛烈的内心震撼。

我想说的是:触发我离开的那桩“特殊诱因”,肯定不是没有\textbf{意义}。至少在我离开后的头几个
月,多半甚至头一年里,这个意义强烈地在场。后来,在新印象的涌入之中、在这场头一波汹涌更新
的动势里,这个意义自然地退居幕后,终于从我眼前消失。可是,即使我不再感知我过去或当下的行
为及其果实的某个“意义”,这个意义并不因此消失。而我重返数学活动——它不可避免地带来同我
离弃的那个世界的更多接触——又冷不丁把我带回那段被遗忘的过去。因为这次“回归”(同当年那
次离开一样出人意料……)最初的收获之一,就是在我那个曾经属于自己的世界里,发现了另
一场腐败——一种据我所知它从未有过的腐败。若要给这件新东西起个名字,涌到我嘴边的是:\textbf{敬
意的丧失}。过去几年里,我不止一次痛切地感到它:当我看到“我昔日所爱的某人,不动声色地踩
压我如今所爱的另一个人——在那人身上,他本该认出我”。在对\emph{葬礼}的思考里,我又一次次重
新遇到它,而且语气更凶,这一回是冲着我亲手生下的某些东西,或者冲着某个胆敢从中汲取灵感的
继承者。在这些时刻,我真正领教了这种丢了敬意的精神的“气息”与“气味”。可我也明白,这种精
神“并不只在我家四邻吹拂”——尽管,恰恰是它吹到我身上、吹到我所爱的人身上,我才真正“认
识”它——正如人们只有亲口尝过,才知道苦果的味道。今天,这种精神已经成了时代精神……

我看得很清楚:这两场腐败——触发我离开的那场,与我“归来”时等着我的那场——并非没有关
联。若要用言词勾勒这层关联的隐隐感受,我要说:\dmark{}\origpage{1223}在科学家们面对军方的金钱(只提这一面)及
其便利时的那种顺手就沾的态度里,我察觉到一种对自己的不敬——个人层面与集体层面皆然
\fnote{1016}{(*)}{在这里,我冒着得罪某些昔日朋友的风险——他们把这种“顺手就沾的态度”当成自己的,却当然不觉得自己亏欠了对自己
的敬意!再说,也绝非确定:别的时代的科学家们,倘若集体地面对同量级的“诱惑”,反应就会不同。常言道,机会造就窃贼!}。而正是在对自己的敬意的丧失里,我认出了对他人的敬意、对出自
他人之手或造物之手的作品的敬意,所以丧失的根由。

我并不声称“理解”了这两场腐败中的任何一场。一方面,有“时代精神”——它的独特动势,在我
看来几乎完全脱出个人行动所能及。这股集体动势,对我始终是个彻底的谜,我从未想过要去探测
它。另一方面,是每一个具体的存在——秉着他的知觉与创造之能,驮着他那全部特殊条件作用的重
负——回应这时代精神的方式;他(自觉或不自觉地)把这种回应,做成了他个人历险的关键要素之
一。

在整个思考过程中,我花了很久,试图去把握\emph{葬礼}的两大主角各自的选择,以及这些选择背
后的力:死者,与葬礼主祭\fnote{1017}{(**)}{(6月22日)到了“最后一刻”,在“家庭相册”一注(n\textsuperscript{o}~173)的 c.(“众人之上的那一位——或首肯”)、d. 与 e. 三部分里,第三个“主角”终于向我显形。}。确定无疑的是,一路走来我学到了东
西;但绝不能说我完成了任务。我甚至可以说,就我的主角而言,我肯定\textbf{没有}完全成功。我收集
拼图的碎片,把它们拼了起来,我甚至相信碎片是对的、拼法大体不错——但对\textbf{整体}的认识,
始终缺着。它至今还是一堆碎片,而这些碎片,对我而言仍是\textbf{陌生的}——陌生于我这个人、我
的经历,因此也就是不可理解的。已做的功课,无疑会有助于我在别的场合多少认路,在该留神的地
方留神(年纪越大,我越发觉要留神的地方真多……)。但这一切,离真正的理解还差得远。于是我
心里冒出一个问题:往这个方向下的功夫,最终是不是一场障眼戏——或者至少,那个\textbf{目标}(在
某种冲突情境里“理解他人”)是不是\dmark{}\origpage{1224}一场障眼戏(虽然走过的\textbf{路}倒是教益丰富……)。我对自己说:真正理解\textbf{那个人}身上的冲突(或者任何一个曾与我亲密相系、而我看
到同样矛盾炸开的人身上的冲突),大概同时也就是\textbf{理解冲突本身}。而且我深知,这样的理解,
不可能来自对他人的冥想(他人永远逃脱我直接的认知),只能来自对我自身的冥想。“阴阳之钥”
这场长考若要显出它的丰饶,靠的不是对他人其人的偶然瞥视,而恰恰是回头看我自己的生命、我
自己的经历,以及我当年对它们的理解。

\subsection{馈赠}   % 18.7.2. Le don

\notehead{180}{}

(4月3日)说到底,我并不觉得有兴致,用几行或几页文字,为\emph{葬礼}中我的主角向我显形的一
切,作一番回顾。照现在的情形,这恐怕不过是一场文体的练习,而对那最残缺不全的理解的更新,
并无多少助益。眼下,我倒是急于抵达这场关于\emph{葬礼}的思考的句号!

当然,我也知道,这个句号绝不等于\emph{葬礼}本身的终结。接下来的几个月,随着这些孤独的果
实般的笔记招来形形色色的回声,想必会有孤独的思考给不了的惊喜与教益。也说不定,来到我这里
的惊喜并不都是苦味的;也许不远的将来还给我备着某种欢乐——正因为它稀罕,才愈发珍贵;正
如仅仅去年(丰收之年!),我就得到过欢乐:收到旧日同事与朋友中的三位——我素来特别敬重或
亲近的——写来的满纸热忱的信\fnote{1018}{(*)}{这里说的是 D.~Mumford、I.M.~Gelfand与 J.~Murre的来信。}。

至于《收获与播种》对数学世界的“时代精神”会有什么全局效应——无论多么微末——几乎不必
说,我对这一点没有半点幻想。至多,这套笔记的出版,或许能叫某桩空前的横逆收场,叫某桩太过刺
眼的不正常得以纠正——而且,我这么说也许还是乐观了。也可能,死者本人的不期复现——他被当
作死了多少年、都馊了——会给那团裹着他双手之作的、压低嗓门的嘲弄合唱,添一个\dmark{}\origpage{1225}终止,或至少
添一个更知趣的弱音器。而倘若这次复现不同时终结那场对一种视野、对一些强而丰
饶的观念的时髦抵制,那它或许至少能打动某一位比别人更大度的年轻数学家,让他毫无保留地从中
汲取灵感(不惜惹人嫌弃),怀着敬意把它们化为己有。

然而,我写《收获与播种》,不是为了这些事里的任何一件——它们当中有的兴许会附带着来,谁
知道呢!我写它,是“为我自己”的,当然,如同我写的一切——作为一场摸索着寻找自身的理解的
手段。但与此同时,对他人的念想——对那些我爱过、而我的历险把我带往\textbf{别处}的那天撇下的
人的念想——贯穿了《收获与播种》写作的始终\fnote{1019}{(*)}{这一念想,在《愚妄与更新》(《收获与播种》的第一部分)里不止一次得到表达。它兴许在后文各部分里不那么显眼,却同样在场。}。这些笔记,既是一场思考,有时也是一场冥
想,对我而言,曾经并且始终是一份\textbf{馈赠}——在我自己之外,送给那些我向他说话的人。当然
我也知道,这份馈赠,也许除了我自己,无人会接。但我并不因此后悔写了它。再说,今天若没有被
某个它属意的人接住,也许明天会被接住。这份既自发、又经长久酝酿的见证,其中每一页、每一个
字都来在其时、来在其所——它明天不会比今天更不真。而无论今天还是明天,若有出乎意料的、
被欣喜迎接的事,那将是得悉:我的馈赠被接住了,哪怕只被一个人——他在我之中认出了他自己……

\subsection{(3) 信使(2)}   % 18.7.3

\notehead{181}{}

比起\emph{葬礼}图景的“前景”,我更不觉得有兴致,就另两个“平面”——一个由“我那班生意兴
隆的学生组成的、扛着铁锹绳索的队伍”,一个由“教团全体”组成——对我在他们身上的所悟与所
惑,作一番原原本本的回顾。关于后者,以及它在\emph{葬礼}里的角色,我在注“掘墓人——或教团全
体”(n\textsuperscript{o}~97)里已经说得相当周详\fnote{1020}{(**)}{(6月22日)我对教团尚且模糊的知觉,最近在已引的注“家庭相册”(n\textsuperscript{o}~173)的 c.、d.、e. 三部分里,出人意料地变得具体了。}。至于我对我那些亲爱的前学生的角色与动机的困惑,它们最清楚地显
现在注“沉默”(n\textsuperscript{o}~84)里,却\dmark{}\origpage{1226}没有在思考的任何后续时刻被认真重审。所以,恰恰
是在这个层次——\emph{葬礼}图景的“第二平面”——我的功课最欠火候!\fnote{1021}{(*)}{(6月22日)不过,关于图景“第二平面”的思考,确有一份(微薄的)续篇:6月19日的注“一场屠杀中的五篇论文——或孝道”(n\textsuperscript{o}~176$_{7}$)。}。在那里,没有一桩功课可与我在
“掘墓人……”那条注里做的相比。图景的这一部分,后来在两道更晚的注里、在阴阳动势的光照下
得到了深化:“天赐的机缘——或封神”与“拒认(一)——或召回”(n\textsuperscript{o}s~151、152)。

“掘墓人——或教团全体”这条注,是思考“第一口气”里写下的最后一道,想必也是它的顶点。可
是,时隔将近一年,我已经不再相信:在我这卑不足道之人的\emph{葬礼}背后(它被看作“对异端的
报复”),那股显得相当明显的集体动机,在集体意志的层面上,真正触到了\emph{葬礼}的\textbf{神
经}。使我生疑的是:在我每一个学生的情形里,这股动机在我看来完全缺席,或者,同其他在场的
力量相比,无关痛痒\fnote{1022}{(**)}{这一事实,出现在注“天鹅绒的脚掌——或微笑”(n\textsuperscript{o}~137)的思考里,打字稿第644--645页。}。可是,\emph{葬礼}里最触目的事实之一,恰恰是它的三个前后相续的“平面”之间
存在的那种“全体一致”:它们的作为与不作为互相衔接、互相补足(仿佛由一个“天衣无缝”的共
同意志指挥着),同一场名副其实的葬礼上一般无二!在如此显著的全体一致里,在心境与作为的如
此整齐划一里,人们当能猜到一个共同的动机,一根同一的“神经”,在驱动着这一个与那一个。

我无意暗示:那股我这里那里察觉到的“弥漫的怨怼”——它起于我的“异端”,(浅浅地)被感受
为开小差,(更深地)被感受为不可容忍的挑衅——全然子虚乌有、全然不起作用。但我如今怀疑,
它的作用是否是决定性的;怀疑它是否就是那根共同的“神经”——难道它人人皆有,却偏偏不
在那些\textbf{在葬礼里戏份最重的人}身上!也就是说,不在我那些学生身上——他们恰恰是那
份遗产的第一批受托人……

这桩“缘由”(我那场“异端”——表面看还相当讲理)在我看来,同那股暴戾之气完全不成比
例——那股气,\dmark{}\origpage{1227}我是在“屠杀一个‘光辉讨论班’”那样的行动里、在教团默许的目光下感到的;同反常层会议上那场当众鼓掌的、同样凶暴的不义,也不成比例。这也不是说,我曾是一
个可憎的同事或老板,怕我怕到积攒的敌意不敢发作、非得等我死了葬了才一齐倾泻到我头上、倾泻
到人们多少“认出我”的那些人头上。从我这里那里收到的种种回声看,既没有\textbf{畏惧}的痕迹
——仿佛我这人曾叫人害怕,如今报了迟到的一箭之仇\fnote{1023}{(*)}{的确,在“愚妄与更新”里,我长篇谈过那种“知名人士”身边——从一个我也说不清何时开始的时刻——笼罩着的\textbf{畏惧},我自己也时或在我周围察觉它的迹象。但那是附着于“知名度”本身的弥散的畏惧,而不是对我这个人本身的——只要稍能建立一点个人接触,它就消失了。我印象里,在个人接触的层面上,人们与其说怕我,不如说把我看作“好脾气”。我确信,在那位“失手——或二十年之后”一节(n\textsuperscript{o}~27)里谈到的学生那里,情形并无不同:每次见面,某种“发怵”持续了很久。这种发怵,今天我看作一种弥漫性的内在不安(“Unsicherheit”)的征兆,它后来在对支配与轻蔑的姿态里找到了补偿与出口。在他众多的学生当中,我有机会认识的三个,个个都结结实实吃过他那副貌似“无缘无故”的恶意姿态的苦头。显然,那种安营扎寨、统治着数学圈子各处的精神,助长了这类悖谬的行径;而这些行径又回过头来塑造这种精神,给它打上那种叫人错愕的、闷声的野蛮印记……};也没有什么稍微具体的行迹可以
归\textbf{罪}于我、足以滋养某种敌意或暴戾(而那种敌意,却从来不报姓名)。

这正是我称作“无缘无故”或“无因”的暴力的典型情境。这种暴力,最终来到“阴阳之钥”这场长
考(它本身构成《收获与播种》的心脏)关注的中心,想来绝非偶然。这种暴力,我认识它远非昨
日始;我头一次同它打照面,并不在我的数学生涯里。即便我有时忘了它在人间世界的存在,那也
从来长不了:它自己会很快前来提醒我,好意不忘。至于当下——借着一个奇怪的“巧合”(我
承认,它常常不合时宜,或至少不受欢迎……),我不记得我此生曾在什么时候,同这样一种暴力的
\dmark{}\origpage{1228}熟悉的征兆打过这么执拗、这么重复、这么磨人的交道——自从我“重返数学”以来,尤其自写《收
获与播种》以来;而最近的这几个月、这几个星期,愈发如此。

\pmark 当然,这里有一个执拗的讯息,它一遍遍回来,而且想必会一直回来,直到被听见。在阴阳长考
的最后几个星期里,我开始侧耳倾听它——同时知道,它要对我说的,我还没有听完。此后两个月
里,一种地下的功课想必在无声中进行。我觉得,那\textbf{要紧而隐秘}\fnote{1024}{(*)}{写下这一行时,我心里清楚,“隐秘”一词在这里是个权宜之计,是对“共识”的让步。我常常有机会发现:当我得知一件我一生都不知道的事时,那件事绝非“隐秘”,恰恰相反,它明摆着、显而易见,有时甚至刺眼——只是我不肯看它而已。发现新事物的情形,大抵如此,无论那是数学的工作,还是认识自己的工作。这种失明的根由、这种常识或直觉官能的阻滞,绝不在这些官能的缺陷,而在于精神的一种近乎不可克服的惰性:它不肯离开共识已然铺好的车辙——无论那是全社会接纳的共识,还是自己所处的那个较小的圈子里的共识,甚至,它们只在我们内心缔结并加盖了印章,如同某位“老板”为他一己之便、同他自己缔结的条约条款……}的东西,开始从表面可见的(或
至少不那么难于承认的)杂事里沉淀出来。“侏儒与巨人”的形象(我朋友皮埃尔供认的)继续萦绕
着我。在这形象背后,我相信辨出一股力量可观的原型:它像童年所受压抑投下的影子,或者
影子之一。它的角色,该是为创造之力所受的压抑提供一个出口与补偿——那压抑,早已内化为那句
“说不出口的无能之信”……在这个预感到的原型里,我仿佛触到一股强大的引擎,驱动着
那些无端施暴的行径,打击那个被当作“巨人”、被当作完好无损之力的携带者的人——这些行径的
爆发,不需要别的“缘由”,只需要一个\textbf{合宜的时机},当下手的风险显得为零、或微不足道。

也许我已经说得太多,而这几行字,才刚刚拂过一道纤细而执拗的直觉——它向我示意:有一桩功
课要做,它还横在我的前头。对这桩功课而言,\emph{葬礼}只是材料之一,同许多别的材料一道,来自
我那所谓“私人”的生活\dmark{}\origpage{1229}。这里不是继续它、甚至不是接近它的地方。它的位置,不在这些
注定要发表的笔记里。

\subsection{(4) 失乐园(2)}   % 18.7.4

\notehead{182}{}

(4月4日)在这场已许诺的回顾——我的思考教给我关于他人的事——里,我的思绪不由自主地
执拗地回到我自己的身上。这对我是个好兆头——它标志着我心中那股回到要紧之事的强烈需要。
对我的认识,额外地带来对他人的理解,而不是相反。自我开始冥想以来,不止一次,“理解他人”
的心思,成了那桩要紧功课——认识我自己——的消遣与岔路。

在有意地回到我自己之前(逆着我那股要抵达那个有名“句号”的不耐!),我还想收入一份最近到
我手里的、关于我朋友皮埃尔的见证。这是我离开数学舞台以来,所知的此类见证中的孤例。它给我
朋友的照明,与我通过别的途径所知的,大异其趣。它也十分适时地再次提醒我:现实,永远比我步
履蹒跚地试图给它画的像更复杂、更丰富\fnote{1025}{(*)}{我绝无意借此暗示:人们(包括我自己不停地)为现实画一幅尽可能“忠实”的像、并随各种“信息”的到来不断调整这幅像的努力,是徒劳而无效的。恰恰相反,这里有一种效力极大的辩证法,让我们同现实接触、去“认识”现实。只有当这幅像(它按事物的本性,自带一份惰性)完全僵死、板结之时,它才成为把握现实的障碍,或者更确切地说,成为让我们的把握官能\textbf{失效}、把我们确实拥有的对现实的知识“蒸发”掉的\textbf{有效手段}。}。

这份见证不是直接的。它关乎一位外国数学家同Deligne的一次(或多或少偶然的)会面留下的印象;
那位同事把这次会面对我的通信人讲了(我猜还是热乎的),通信人又在一封信里把经过转述给我。
经我的通信人与那位讲述会面的同事(下文称他“Z”)同意,我在这里翻译信中涉及这次会面的部分。
我的通信人推测,这一幕应当发生在1981年。(注意:这也正是反常层会议之年——不过通信人与我
之间,倒不曾谈起过那场讨论会。)

\begin{quote}
\pmark\origpage{1230}“……有一天,Z.~去 Bures 参加一个会,发现自己置身一间上茶的屋子[看
来是法国高等科学研究所的‘茶室’],屋里有很多数学家。只见门开了,Deligne走进来。Z.~先生把
这一幕讲得相当生动:他看上去蔫蔫的,双臂低垂,人觉出他周身有一种孤立。别的人都定定地看
着他,有点像看一只珍禽,又没人知道能对他说点什么。Z.~靠窗独自坐着,Deligne有些迟疑,在他
身旁坐了下来。Z.~不知道跟他说什么好。这时他起了个念头,就简单地说:他觉得围绕‘平展拓扑’
等等的那一整套想法,还有您带来的那些新思想,非同寻常。[“您”,这里及下文,指的都是我,Grothendieck——我的通信人是转述给我的。]Deligne的眼睛顿时亮了,他对 Z.~说:是啊,那是数学里
最好的东西之一;听您的报告\fnote{1026}{(*)}{同前,这里的“您”指的就是我,Grothendieck。},那才叫美……他接着说:您想想这个,再想想那个……数出一大
串事情来——Z.~自己(据他向我说的)一窍不通,但他看见了谈话对方身上骤然涌起的热情。Deligne又说:多可惜,您\textsuperscript{(*)}退了!他确信,晶体上同调和许多别的东西,不至于落到现在这副半
生不熟的样子;要是您\textsuperscript{(*)}真还接着干,它们如今早就是站得稳稳的构造了,就像平展上同调
那样……”
\end{quote}

这个故事里有两件事击中了我。其一是那种孤立的印象——它似乎深深击中了 Z.~先生。这一印象究
竟出自Deligne生涯里某个特别的时刻,还是这样的孤立终于浸透了他同所有同侪的关系,我很难说。
我没有收到过任何指向后者的别的见证。

另一件同样触目的事,在回到我这里的种种回声之中同样独一无二,就是提到我的名字与某段过去
时,那份热情的骤然涌现。那段过去,他早已决意宣告无效作废。还有他在这段过去里的根。而在那段
过去里,也还有一种童年的清新——那种被他逐出自己的“成人”生活、\dmark{}\origpage{1231}逐出这个重要而受人仰慕的
男子的生活的清新。在他周围,识相的人想必都不去提那段过去,不去提他还只是众多学生
中的一个、还怀着一腔美好激情的那些年月——正如在阔人宅邸里,满屋老样式的家具,谁也不去
提寒素甚至窘迫的起家……

而偏偏是这个陌生人,鬼使神差地坐在他身旁,忽然热切地谈起这件从来无人谈起的事(至少不当着
他的面谈……)!当然,这就像那团精选的、拘谨的空气骤然消散了;一个陌生人的热情,唤醒了他
身上同样的热情,并且——在一瞬之间——把他重新接通到那眼遥远的泉源,那眼他被认为早已忘
却、早已失去的泉源……
